\subsection{Comparação das previsões com referências observáveis}

A avaliação complementar utiliza os mesmos artefatos quantitativos e os horizontes escolhidos em 2024. O alvo é o log-retorno acumulado em 5, 10 ou 21 pregões. Para cada ativo e horizonte, o modelo e as referências são avaliados nos mesmos pares, com desfechos terminando até 30 de dezembro de 2025. Não houve retreinamento ou escolha de novos modelos usando os resultados abaixo.

São utilizadas três referências: retorno zero; média expansiva dos log-retornos de mesmo horizonte já realizados; e direção mais frequente entre esses alvos históricos. As duas referências históricas exigem pelo menos 80 alvos conhecidos. Na previsão ao fechamento de t, o último alvo admitido é aquele iniciado em t-h e concluído em t. A direção majoritária é somente uma regra de classificação; não lhe é atribuída magnitude ou RMSE. Diferentemente da maior classe observada em 2025, apresentada anteriormente como diagnóstico ex post, esta referência é calculável sem consultar desfechos futuros.

Define-se $S_0=1-\mathrm{MSE}_{XGB}/\mathrm{MSE}_{0}$. Valores positivos representam melhora contra o retorno zero. O indicador $S_{hist}$ substitui o denominador pelo MSE da média histórica. Essa comparação respeita a escala de cada ativo e horizonte; RMSEs brutos de horizontes distintos não são comparados como problemas equivalentes.

\begin{table}[H]\centering
\caption{Qualidade preditiva dos horizontes selecionados em 2024.}
{\footnotesize\begin{tabular}{lrrrrrr}\hline
Ativo & h & HR XGB & HR maj. & HR média & $100S_0$ & $100S_{hist}$\\ \hline
ABEV3 & 5 & 45,71 & 48,57 & 48,57 & -6,07 & -5,05\\
B3SA3 & 21 & 53,28 & 61,14 & 62,88 & -2,14 & -6,51\\
BBAS3 & 10 & 43,75 & 50,42 & 50,42 & -13,56 & -11,55\\
BBDC4 & 5 & 52,65 & 60,41 & 53,47 & 0,87 & 0,77\\
CMIG4 & 21 & 61,14 & 65,94 & 65,94 & -24,00 & -33,14\\
EQTL3 & 5 & 54,69 & 60,82 & 60,82 & 0,02 & -5,06\\
GGBR4 & 21 & 44,10 & 61,57 & 61,57 & -23,26 & -28,15\\
ITUB4 & 21 & 65,07 & 76,86 & 76,86 & 3,16 & -8,98\\
LREN3 & 10 & 47,08 & 40,83 & 40,83 & -1,06 & 0,08\\
MGLU3 & 21 & 47,60 & 52,40 & 47,60 & -37,10 & -31,91\\
PETR3 & 10 & 46,67 & 46,67 & 46,67 & -10,48 & -3,82\\
PETR4 & 21 & 54,59 & 48,47 & 48,47 & -8,39 & 10,49\\
PRIO3 & 21 & 48,47 & 48,47 & 48,47 & -20,34 & 3,64\\
RADL3 & 21 & 67,25 & 58,52 & 58,08 & -1,08 & -0,37\\
RAIL3 & 21 & 55,46 & 44,98 & 44,98 & -0,17 & 1,60\\
SANB11 & 21 & 56,33 & 57,21 & 64,19 & 2,77 & 0,10\\
SBSP3 & 21 & 70,31 & 70,31 & 70,31 & 29,52 & 12,34\\
SUZB3 & 10 & 37,08 & 37,08 & 37,08 & -6,20 & 4,75\\
VALE3 & 10 & 33,75 & 71,25 & 71,25 & -10,99 & -20,32\\
WEGE3 & 21 & 63,76 & 40,61 & 40,61 & 16,13 & 25,04\\
\hline\end{tabular}}\par\smallskip{\footnotesize Fonte: elaboração própria.}\end{table}

A acurácia média entre ativos é 52,44\%, frente a 55,13\% para a direção majoritária histórica e 54,95\% para o sinal da média histórica. O XGBoost supera estritamente a direção majoritária em 5 ativos e empata em 4. Em magnitude, melhora o RMSE do retorno zero em 6 ativos e o da média histórica em 9. Os resultados distinguem capacidade de acertar o sinal e capacidade de estimar a magnitude. Valores de HR e skills na tabela estão em porcentagem.

LREN3, PETR4, RADL3, RAIL3 e WEGE3 superam a direção majoritária histórica. WEGE3 também reduz o erro frente às duas referências de magnitude. SBSP3 apresenta elevada acurácia, mas empata com a direção majoritária; isso mostra por que um percentual alto de acertos, isoladamente, não certifica ganho preditivo. O skill frente ao retorno zero e o R-quadrado usam referências diferentes: é possível superar o retorno zero e ainda apresentar R-quadrado negativo frente à média retrospectiva do teste.

\subsection{Horizontes e estabilidade temporal}

\begin{table}[H]\centering
\caption{Diagnóstico de todos os horizontes em 2025, sem nova seleção.}
{\footnotesize\begin{tabular}{rrrrr}\hline
h & HR XGB (\%) & HR maj. (\%) & $100S_0$ médio & Ativos com melhora\\ \hline
5 & 50,67 & 51,78 & -3,81 & 6\\
10 & 50,27 & 54,31 & -2,55 & 6\\
21 & 54,65 & 54,83 & -10,33 & 5\\
\hline\end{tabular}}\par\smallskip{\footnotesize Fonte: elaboração própria.}\end{table}

Os três horizontes são mostrados como diagnóstico retrospectivo. Seu desempenho em 2025 não é utilizado para substituir a escolha feita em 2024. Cada horizonte possui quantidade diferente de alvos elegíveis e graus distintos de sobreposição.

\begin{figure}[H]\centering
\caption{Heterogeneidade do erro preditivo por ativo e horizonte.}
\includegraphics[width=0.92\textwidth]{05_Resultados/predicao_xgb/skill_horizontes.png}
\par\smallskip{\footnotesize Fonte: elaboração própria.}\end{figure}

\begin{table}[H]\centering
\caption{Estabilidade trimestral dos horizontes previamente selecionados.}
{\footnotesize\begin{tabular}{lrrr}\hline
Trimestre & HR XGB (\%) & HR maj. (\%) & HR média (\%)\\ \hline
2025Q1 & 52,87 & 52,21 & 50,25\\
2025Q2 & 55,90 & 50,57 & 53,20\\
2025Q3 & 50,30 & 54,02 & 53,41\\
2025Q4 & 50,95 & 67,78 & 66,68\\
\hline\end{tabular}}\par\smallskip{\footnotesize Fonte: elaboração própria.}\end{table}

\begin{figure}[H]\centering
\caption{Comparação trimestral com referências disponíveis no instante da previsão.}
\includegraphics[width=0.92\textwidth]{05_Resultados/predicao_xgb/estabilidade_trimestral.png}
\par\smallskip{\footnotesize Fonte: elaboração própria.}\end{figure}

No segundo trimestre, a acurácia média do XGBoost é 55,90\%, contra 50,57\% da direção majoritária histórica. No quarto trimestre, o modelo alcança 50,95\%, enquanto essa referência atinge 67,78\%. Essa mudança evidencia dependência temporal do desempenho, sem permitir identificar sua causa somente pelas métricas agregadas.

As médias trimestrais atribuem o mesmo peso a cada ativo. O trimestre corresponde à data de origem da previsão; um alvo pode terminar no trimestre seguinte, desde que concluído até o final do teste. Os alvos sobrepostos e os choques comuns entre ações impedem tratar os pares como observações independentes. As diferenças apresentadas são descritivas, sem interpretação automática de significância estatística.

\subsection{Atributos e limites da explicabilidade}

\begin{figure}[H]\centering
\caption{Participação relativa dos atributos nos resumos SHAP históricos.}
\includegraphics[width=0.92\textwidth]{05_Resultados/predicao_xgb/shap_historico.png}
\par\smallskip{\footnotesize Fonte: elaboração própria.}\end{figure}

Nos resumos disponíveis, os atributos de momentum de 252 e 126 pregões possuem as maiores participações médias, seguidos por medidas de risco de mercado, volatilidade e tendência de mercado de 63 pregões. A presença dessas variáveis descreve os modelos salvos; não prova que elas tenham acrescentado poder preditivo fora da amostra.

Para cada ativo, normaliza-se a média do valor absoluto de SHAP pela soma de seus atributos; as participações são então promediadas igualmente entre os vinte ativos. Valores absolutos informam intensidade de contribuição, sem indicar se um atributo eleva ou reduz a previsão. Os arquivos disponíveis são resumos dos modelos salvos, e não explicações individuais das previsões OOF de 2025. Portanto, a figura não demonstra estabilidade temporal, causalidade econômica ou importância exclusiva no teste.

O conjunto de evidências aponta desempenho heterogêneo e capacidade limitada de previsão de magnitude. O retorno da carteira não constitui validação isolada do XGBoost. Embora o horizonte tenha sido escolhido em 2024, os artefatos preservam a seleção original entre regressor e classificador, que pode consultar o OOF completo. A validação interna de hiperparâmetros também precisa ser auditada quanto à separação de alvos sobrepostos. Uma avaliação confirmatória exige reconstruir treino e seleção de forma aninhada antes de um novo período de teste.
